% =============================================================================
% Chapter 1: Overview
% =============================================================================
\chapter{Overview}
\label{ch:overview}

\section{Purpose of This Document}

This document serves as the comprehensive verification report and usage guide for the
\textbf{APB SPI Master IP} (\texttt{apb\_spi\_master}), a configurable Serial Peripheral Interface
(SPI) master controller with an Advanced Peripheral Bus (APB) register interface. The IP originates
from the PULP (Parallel Ultra-Low-Power Processing) platform developed by ETH Zurich and the
University of Bologna, licensed under the Solderpad Hardware License v0.51.

This report is intended for:
\begin{itemize}[leftmargin=2em]
  \item \textbf{Verification engineers} who need to understand how the testbench works and what
        each testcase proves.
  \item \textbf{Design engineers} who need to integrate the IP into a System-on-Chip (SoC) and
        configure it correctly.
  \item \textbf{New team members} who need a self-contained reference to safely use and extend the IP.
  \item \textbf{Teams working on other APB peripherals} (Timer, GPIO, UART) who can reuse the
        testbench methodology described here.
\end{itemize}

\section{What the IP Is and What Problem It Solves}

The APB SPI Master IP provides a processor-accessible SPI master controller. In an embedded SoC,
processors communicate with external devices---flash memories, sensors, ADCs, DACs, and other
peripherals---using the SPI protocol. Directly bit-banging SPI from software is slow and wastes
CPU cycles. A dedicated hardware SPI master offloads this work: the processor writes configuration
registers and data via the APB bus, and the hardware autonomously generates the SPI clock, drives
chip-select lines, serializes transmit data, and deserializes receive data.

Key capabilities of this IP:
\begin{enumerate}[leftmargin=2em]
  \item \textbf{Standard and Quad SPI:} Supports single-bit (standard) and four-bit (quad) data
        transfer modes for higher throughput.
  \item \textbf{Configurable clock divider:} An 8-bit programmable divider allows SPI clock
        frequencies from $f_{sys}/2$ down to $f_{sys}/512$.
  \item \textbf{Flexible frame format:} Independent command, address, dummy, and data phases with
        programmable bit lengths up to 16 bits for command/address and 65,535 bits for data.
  \item \textbf{Four chip-select outputs:} A 4-bit chip-select register allows addressing up to
        four SPI slave devices without external decode logic.
  \item \textbf{Deep FIFOs:} Parameterizable TX and RX FIFOs (default depth 10) decouple the
        processor from the SPI transfer timing.
  \item \textbf{Interrupt support:} Configurable threshold-based and counter-based interrupts for
        TX watermark and RX watermark events, plus an End-of-Transfer (EOT) event output.
  \item \textbf{Software reset:} A register-accessible reset clears the FIFOs without requiring
        a full system reset.
\end{enumerate}

\section{Scope of the Verification Effort}

The verification effort aimed to prove the following:
\begin{enumerate}[leftmargin=2em]
  \item The IP resets to a known, safe idle state.
  \item All configuration registers are writable and readable with correct bit-field mapping.
  \item Standard SPI write and read transactions complete correctly with proper EOT signaling.
  \item The receive data path correctly captures data from the SPI bus and delivers it through the
        RX FIFO.
  \item Complex multi-phase transactions (command + address + dummy + data) function correctly.
  \item Quad SPI write and read modes operate correctly.
  \item The TX FIFO supports burst loading and sequential transmission.
  \item Multi-word RX transfers fill the RX FIFO with multiple 32-bit words.
  \item The chip-select multiplexer correctly selects alternate slave devices.
  \item Software reset clears both FIFOs and returns the controller to idle.
  \item The EOT event output produces a single-cycle pulse.
  \item No spurious interrupts occur after reset and interrupt enable.
\end{enumerate}

The testbench (\texttt{tb\_apb\_spi\_master} v2) implements 14 testcases (TC1--TC14) containing
28 individual pass/fail checks. All 28 checks passed in the final Xcelium run, completing at
simulation time 27,435~ns.

\section{Acronyms and Abbreviations}

\begin{longtable}{@{}ll@{}}
\toprule
\textbf{Acronym} & \textbf{Definition} \\
\midrule
\endhead
APB     & Advanced Peripheral Bus (AMBA 3) \\
BFM     & Bus Functional Model \\
CMD     & Command phase \\
CS / CSN & Chip Select (active-low) \\
DAC     & Digital-to-Analog Converter \\
DUT     & Device Under Test \\
EOT     & End of Transfer \\
FIFO    & First-In, First-Out buffer \\
FSM     & Finite State Machine \\
GPIO    & General-Purpose Input/Output \\
IP      & Intellectual Property (reusable hardware block) \\
LSB     & Least Significant Bit \\
MSB     & Most Significant Bit \\
MOSI    & Master Out Slave In (SDO in this IP) \\
MISO    & Master In Slave Out (SDI in this IP) \\
PRDATA  & APB read data bus \\
PULP    & Parallel Ultra-Low-Power Processing platform \\
PWDATA  & APB write data bus \\
QSPI    & Quad SPI (4-bit data bus) \\
RTL     & Register Transfer Level \\
RX      & Receive \\
SDI     & Serial Data In (from slave to master) \\
SDO     & Serial Data Out (from master to slave) \\
SoC     & System on Chip \\
SPI     & Serial Peripheral Interface \\
TX      & Transmit \\
UART    & Universal Asynchronous Receiver/Transmitter \\
\bottomrule
\end{longtable}
